\section{The ten-dimensional algebra}
\label{a:section}

We now construct an algebra \(C\) and a simple module \(s\) whose
resolution supplies the degree-three extension classes used below.
Besides this resolution, we need
finite projective-dimension bounds for the dual regular modules;
these bounds make derived induction admit a finite bimodule
representative. The multiplication of $C$ is the one in
\cite[Section~4.1]{OpenAIAR2026}; all properties needed here are proved
directly.  We work over the fixed field
\(k=\mathbb F_2(q,H_1,H_2)\), where the three parameters are
algebraically independent.

Let \(C\) have the following basis, arranged by its left and right
idempotents:
\begin{equation}
\begin{array}{c|c@{\qquad}c|c}
\text{corner}&\text{basis}&\text{corner}&\text{basis}\\ \hline
eCe&e,x,y,z&eCf&u,v\\
fCe&t,j&fCf&f,n.
\end{array}
\label{a:corners}
\end{equation}
The unit is \(1=e+f\).  The complete multiplication table, with the
left factor indexing the row, is
\begin{equation}
\begingroup
\setlength{\arraycolsep}{4pt}
\begin{array}{c|cccccccccc}
 &e&f&x&y&z&u&v&t&j&n\\ \hline
e&e&0&x&y&z&u&v&0&0&0\\
f&0&f&0&0&0&0&0&t&j&n\\
x&x&0&0&qz&0&v&0&0&0&0\\
y&y&0&z&0&0&v&0&0&0&0\\
z&z&0&0&0&0&0&0&0&0&0\\
u&0&u&0&0&0&0&0&y+qx&z&v\\
v&0&v&0&0&0&0&0&qz&0&0\\
t&t&0&j&q^2j&0&(1+q)n&0&0&0&0\\
j&j&0&0&0&0&0&0&0&0&0\\
n&0&n&0&0&0&0&0&qj&0&0
\end{array}
\endgroup
\label{a:multiplication}
\end{equation}

\begin{lemma}\label{a:table}
The table \eqref{a:multiplication} defines an associative unital algebra
of dimension ten.  Its radical is the span of the eight basis vectors
other than \(e,f\), and \((\rad C)^5=0\).
\end{lemma}

\begin{proof}
Give the basis vectors the degrees
\[
\begin{gathered}
\deg e=\deg f=0,\qquad \deg u=\deg t=1,\\
\deg x=\deg y=\deg n=2,\qquad
\deg v=\deg j=3,\qquad \deg z=4.
\end{gathered}
\]
The table is homogeneous and respects the corners.  Thus associativity
with an idempotent, or with incompatible adjacent corners, follows
from \eqref{a:corners}.  For positive-degree triples of total degree
greater than four, both bracketings vanish.  The remaining composable
triples and the value of either bracketing are precisely
\[
\begin{array}{c|ccccc}
abc&utu&tut&xut&yut&ntu\\ \hline
(ab)c=a(bc)&(1+q)v&q(1+q)j&qz&qz&0
\end{array}
\]
\[
\begin{array}{c|cccccc}
abc&unt&txu&tyu&utx&uty&tun\\ \hline
(ab)c=a(bc)&qz&0&0&z&q^2z&0.
\end{array}
\]
Every entry is obtained by two applications of
\eqref{a:multiplication}; this exhausts the associativity check by
trilinearity.  The positive-degree span is a nilpotent ideal with
fifth power zero and quotient \(ke\oplus kf\cong k^2\).  It is therefore
the radical.
\end{proof}

Let \(s\) be the one-dimensional left \(C\)-module on which \(f\)
acts as the identity and every other basis vector acts as zero.
Its vector space dual \(Ds\) is the right simple with the same character.

The corner \(eCe\) gives a concrete connection with Schulz's
example~\cite[Example~7]{Schulz1995}. The table identifies it with
the quantum exterior algebra
\[
 eCe\cong k\langle x,y\rangle/(x^2,y^2,yx-q^{-1}xy).
\]
Moreover, the assignments
\(t\mapsto x+q^{-1}y\) and \(j\mapsto q^{-1}z\) give an
isomorphism of right \(eCe\)-modules
\[
 fCe\cong(x+q^{-1}y)eCe.
\]
The identities \(tx=j\), \(ty=q^2j\), and \(jx=jy=0\) verify
right-linearity, and the two images form a basis of the right ideal.
With Schulz's
parameter \(p=q^{-1}\), this right ideal is the first syzygy of his
module \((x+y)eCe\). The later left-module resolutions will use the
same type of parameter shift within the ten-dimensional algebra.

Write \(a^*\) for the basis vector of \(DC\) dual to a basis vector
\(a\) of \(C\).  Form the trivial extension
\begin{equation}
T=C\ltimes DC,\qquad
(a,\phi)(b,\psi)=(ab,a\psi+\phi b),\qquad
(a\phi b)(c)=\phi(bca).
\label{a:trivial-extension}
\end{equation}
By Lemma~\ref{tr:trivial-symmetric}, this is a symmetric algebra of
dimension twenty, with symmetrizing functional
\(t_T(a,\phi)=\phi(1)\).  In particular,
\[
t_T\bigl((a,\phi)(b,\psi)\bigr)=\phi(b)+\psi(a).
\]
We regard \(s\) as a \(T\)-module through the quotient \(T\to C\).
Set
\begin{equation}
B=C\otimes C\subseteq E=T\otimes T.
\label{a:tensor-algebras}
\end{equation}
Then \(\dim_k B=100\) and \(\dim_k E=400\).  The functional
\(t_E=t_T\otimes t_T\) gives a symmetric nondegenerate product pairing,
so \(E\) is symmetric.

We record the radical and corner information that will enter the trace
argument.  One has
\begin{align}
\rad T&=(\rad C)\oplus DC,\nonumber\\
\rad B&=(\rad C)\otimes C+C\otimes(\rad C),\nonumber\\
\rad E&=(\rad T)\otimes T+T\otimes(\rad T).
\label{a:radicals}
\end{align}
For the first equality, the displayed ideal is a square-zero extension
of the nilpotent ideal \(\rad C\), and its quotient is \(k^2\).
For the tensor products the displayed ideals are sums of commuting
nilpotent ideals, and their quotients are \(k^4\).  This proves all
three equalities.  The four vertex idempotents of \(B\), also used in
\(E\), are \(e\otimes e,e\otimes f,f\otimes e,f\otimes f\).
For these vertices the corner dimensions are
\begin{equation}
\bigl(\dim_k(pB p')\bigr)_{p,p'}=
\begin{pmatrix}
16&8&8&4\\
8&8&4&4\\
8&4&8&4\\
4&4&4&4
\end{pmatrix}.
\label{a:corner-parity}
\end{equation}
Indeed these dimensions are the products of the corresponding corner
dimensions \(4,2,2,2\) of \(C\).  In particular, every entry is even.

For \(i\in\mathbb Z\), put
\[
\ell_i=x+q^i y.
\]
We denote right multiplication by \(a\) by \(\rho_a\) and left
multiplication by \(a\) by \(\lambda_a\).  Thus \(\rho_a\) is a map
of left modules and \(\lambda_a\) is a map of right modules whenever
the indicated source and target corners allow it.

\begin{proposition}\label{a:dual-dimensions}
The dual regular module \(DC\) has projective dimension at most two
on both sides.  More precisely,
\[
D(eC)\cong Ce,\qquad D(Ce)\cong eC,
\]
and there are exact sequences
\begin{align}
0\longrightarrow Cf\xrightarrow{\rho_t}Ce
 \xrightarrow{\rho_{\ell_{-2}}}Ce
 \xrightarrow{e\mapsto j^*}D(fC)\longrightarrow0,
\label{a:left-dual-resolution}\\
0\longrightarrow fC\xrightarrow{\lambda_u}eC
 \xrightarrow{\lambda_{\ell_0}}eC
 \xrightarrow{e\mapsto v^*}D(Cf)\longrightarrow0.
\label{a:right-dual-resolution}
\end{align}
Consequently
\begin{equation}
\pd_B({}_B E)\le4,\qquad
\pd_{B^{\mathrm{op}}}(E_B)\le4,\qquad
\idim_B({}_B B)\le4.
\label{a:dimension-bounds}
\end{equation}
\end{proposition}

\begin{proof}
The pairing \(eC\times Ce\to k\) that takes the coefficient of \(z\)
in a product is nondegenerate.  In the ordered bases
\[
\mathcal B_{eC}=(e,x,y,z,u,v),\qquad
\mathcal B_{Ce}=(e,x,y,z,t,j),
\]
its only nonzero entries are the products
\(ez,xy,yx,ze,uj,vt\), with coefficients \(1,q,1,1,1,q\).
Thus it has determinant \(q^2\ne0\).  The maps sending \(e\) to
\(z^*\) give the stated left- and right-module isomorphisms.

Here is the full calculation for the other dual summands.  The two
surjections, on the ordered bases just specified, have the images
\begin{align*}
Ce\longrightarrow D(fC):\quad
(e,x,y,z,t,j)&\longmapsto
(j^*,t^*,q^2t^*,0,qn^*,f^*),\\
eC\longrightarrow D(Cf):\quad
(e,x,y,z,u,v)&\longmapsto
(v^*,u^*,u^*,0,n^*,f^*).
\end{align*}
Each has rank four; their kernels are respectively
\(\langle\ell_{-2},z\rangle_k\) and
\(\langle\ell_0,z\rangle_k\).
The intermediate maps have images and kernels
\[
\begin{array}{c|c|c}
\text{map}&\text{image}&\text{kernel}\\ \hline
\rho_{\ell_{-2}}:Ce\to Ce
 &\langle\ell_{-2},z\rangle_k
 &\langle y+qx,z,t,j\rangle_k\\
\lambda_{\ell_0}:eC\to eC
 &\langle\ell_0,z\rangle_k
 &\langle y+qx,z,u,v\rangle_k.
\end{array}
\]
For completeness, the initial injections are given by
\begin{align*}
\rho_t:(f,u,v,n)&\longmapsto(t,y+qx,qz,qj),\\
\lambda_u:(f,t,j,n)&\longmapsto(u,y+qx,z,v).
\end{align*}
Their four images are independent and span the respective kernels in
the table.  This proves exactness of
\eqref{a:left-dual-resolution} and
\eqref{a:right-dual-resolution}, including their orientations and
the exceptional index \(-2\) in the first sequence.

On either \(C\)-side, \(T=C\oplus DC\) has a projective resolution of
length at most two.  Tensor two such resolutions over the field and
totalize.  The resulting resolutions of the corresponding
\(B\)-modules \(E\) have length at most four, proving the first two
bounds in \eqref{a:dimension-bounds}.  Likewise
\(DB\cong DC\otimes DC\) has right \(B\)-projective dimension at most
four.  Dualizing a finite projective resolution gives an injective
coresolution of \({}_B B\), and proves the last bound.
\end{proof}

We now resolve \(s\). The same multiplication maps will also determine
the dual complex, including its right-module cohomology.

\begin{proposition}\label{a:resolution}
The augmentation \(Cf\to s\) extends to a projective resolution
\(R\to s\) with
\begin{equation}
R_0=Cf,\qquad R_i=Ce\ (i\ge1),\qquad
d_1=\rho_u,\qquad d_{i+2}=\rho_{\ell_i}\ (i\ge0).
\label{a:simple-resolution}
\end{equation}
It satisfies
\begin{equation}
H^a\Hom_C(R,s)=
\begin{cases}k&a=0,\\0&a\ne0,\end{cases}
\qquad
H^a\Hom_C(R,C)=
\begin{cases}Ds&a=2,\\0&a\ne2,\end{cases}
\label{a:dual-cohomology}
\end{equation}
where the second identification is an isomorphism of right
\(C\)-modules.
\end{proposition}

\begin{proof}
On the basis \((e,x,y,z,t,j)\) of \(Ce\), right multiplication gives
\begin{align*}
\rho_u(e,x,y,z,t,j)
  &=(u,v,v,0,(1+q)n,0),\\
\rho_{\ell_i}(e,x,y,z,t,j)
  &=(\ell_i,q^{i+1}z,z,0,(1+q^{i+2})j,0).
\end{align*}
For every \(i\ge0\), the coefficient \(1+q^{i+2}\) is nonzero.
Consequently
\[
\begin{array}{c|c|c}
\text{map}&\text{image}&\text{kernel}\\ \hline
\rho_u
 &\langle u,v,n\rangle_k
 &\langle\ell_0,z,j\rangle_k\\
\rho_{\ell_i},\ i\ge0
 &\langle\ell_i,z,j\rangle_k
 &\langle\ell_{i+1},z,j\rangle_k.
\end{array}
\]
The first image is \(\rad(Cf)\), and each later image is the preceding
kernel.  This proves exactness at every index.  Since
\(\Hom_C(Ce,s)=es=0\) and \(\Hom_C(Cf,s)=fs=k\), the first equality
in \eqref{a:dual-cohomology} follows.

Under \(\Hom_C(Cf,C)=fC\) and \(\Hom_C(Ce,C)=eC\), the dual cochain
complex is
\begin{equation}
fC\xrightarrow{\lambda_u}eC
 \xrightarrow{\lambda_{\ell_0}}eC
 \xrightarrow{\lambda_{\ell_1}}eC
 \xrightarrow{\lambda_{\ell_2}}\cdots,
\label{a:dual-complex}
\end{equation}
starting in degree zero.  The first arrow is injective, with image
\(\langle u,y+qx,z,v\rangle_k\).  For \(i\ge0\),
\[
\lambda_{\ell_i}(e,x,y,z,u,v)
  =(\ell_i,q^iz,qz,0,(1+q^i)v,0).
\]
The image and kernel are therefore
\[
\begin{array}{c|c|c}
 &\text{image}&\text{kernel}\\ \hline
i=0&\langle\ell_0,z\rangle_k
 &\langle\ell_{-1},z,u,v\rangle_k\\
i\ge1&\langle\ell_i,z,v\rangle_k
 &\langle\ell_{i-1},z,v\rangle_k.
\end{array}
\]
Since \(q\ell_{-1}=qx+y\), the first kernel equals the image of
\(\lambda_u\).  The only remaining cohomology is in degree two:
\[
\frac{\langle\ell_0,z,v\rangle_k}
     {\langle\ell_0,z\rangle_k}=k\overline v.
\]
Finally \(vf=v\), \(ve=0\), and the only nonzero product of \(v\)
with a radical basis vector is \(vt=qz\), which is a boundary.
Thus the radical acts trivially on \(k\overline v\), and its right
module structure is exactly \(Ds\).
\end{proof}
